\section{Alzheimer's：从机制复现到一个前瞻节点}

Alzheimer's 案例比前面的隐藏结果复现更复杂：系统先重建一个多步骤 ACE--bradykinin--B2R 机制链，再由合作团队对其中一个节点做前瞻实验。它因此同时包含公开背景、回顾性 recapitulation 和未发表 prospective validation，必须拆开叙述。

\subsection{先标明合作与发表状态}

课程与 MGH 研究者合作，目标是复现多年人类实验路线，并在数日内独立预测 ACE--B2R 致病轴。slide 同样写明 unpublished data, needs peer review。这里的“case study”不是临床病例，也不是已发表治疗建议，而是机制假设与实验流程展示。

\lecturefigure{slide-016.jpg}{Alzheimer's 案例的目标与未发表边界}{Stanford CS25 V6 Lecture 07 官方录像 00:53:54}

\begin{warningbox}{三类证据必须分栏}
公开文献已经讨论 bradykinin、认知损伤和 Alzheimer’s 的联系；AI co-scientist 重建 ACE inhibitor 到 B2R 的机制链；合作团队对 B2R 节点做了课堂所述前瞻实验。前两者不能证明第三者，第三者在未发表前也不能被写成已确证治疗机制。
\end{warningbox}

\subsection{ACE--bradykinin--B2R 链怎样形成}

ACE inhibitor 用于高血压治疗，会抑制 angiotensin-converting enzyme。ACE 除参与肾素--血管紧张素系统，也降解 bradykinin；抑制 ACE 可能提高 bradykinin 信号。课程提出：bradykinin 激活 B2R 后触发炎症级联，可能解释部分人群中 ACE inhibitor 与 Alzheimer’s 风险的观察关联。

\lecturefigure{slide-017.jpg}{AI 对 ACE inhibitor、bradykinin 与 B2R 机制链的复现}{Stanford CS25 V6 Lecture 07 官方录像 00:54:48}

\begin{knowledgebox}{读图：四步链条需要逐步验证}
左侧列出临床悖论：ACE inhibitor 常用于高血压，但某些观察中与 Alzheimer’s 风险相关；右侧是 AI 构造的机制。应分别问：药物是否改变 bradykinin？bradykinin 是否通过 B2R 作用于相关细胞？B2R 是否导致病理表型？阻断 B2R 是否能逆转？只要一环失败，完整因果叙述就需要修订。
\end{knowledgebox}

可以把机制链写成结构化因果假设：
$$
\text{ACE inhibition}
\rightarrow \uparrow\text{bradykinin}
\rightarrow \text{B2R activation}
\rightarrow \text{inflammatory phenotype}.
$$
其中：
\begin{itemize}
  \item 第一个箭头依赖 ACE 对 bradykinin 的降解作用；
  \item 第二个箭头是配体--受体关系；
  \item 第三个箭头才连接到神经炎症或病理表型；
  \item 观察相关不能自动确定箭头方向，也不能排除共同病因与用药偏差。
\end{itemize}

\begin{warningbox}{Causal story 的流畅性会制造过度确信}
LLM 很擅长把多个真实局部事实串成完整链条，但每条边的证据可能来自不同物种、组织或实验条件。科研 Agent 必须输出 edge-level provenance 和可反驳实验，而不是只给一段顺滑机制叙述。
\end{warningbox}

\subsection{Prospective validation：为什么一个节点很重要}

课堂称系统识别了一个被忽视的 regulatory node，并提出使用 antagonist 干预 B2R。合作团队在实验室中测试预测，报告了新 inhibitor 成功 reverse disease-associated phenotype。若候选是在实验前冻结，这比回顾性解释更强，因为系统对结果承担了可失败的预测风险。

\lecturefigure{slide-018.jpg}{B2R 节点的预测、实验与验证三阶段}{Stanford CS25 V6 Lecture 07 官方录像 00:55:26；课堂未发表数据}

\begin{knowledgebox}{读图：prospective 的关键是时间顺序}
第一列 Prediction 必须发生在看到实验结果之前；第二列 Experiment 把 B2R antagonist 变成可执行干预；第三列 Validation 报告表型逆转。若这三步有时间戳、冻结报告和预先定义指标，就能减少事后挑选。图仍未展示完整方法和独立重复，所以强于 recapitulation，弱于公开可复核证据。
\end{knowledgebox}

若以实验前证据 $D_0$ 形成假设 $h$，实验结果为 $e$，更新可以写成
$$
p(h\mid D_0,e)
=\frac{p(e\mid h,D_0)p(h\mid D_0)}
{p(e\mid D_0)}.
$$
其中：
\begin{itemize}
  \item $p(h\mid D_0)$ 是实验前可信度；
  \item $p(e\mid h,D_0)$ 是若机制成立，出现结果 $e$ 的概率；
  \item $p(e\mid D_0)$ 包含其他机制也产生同一结果的可能；
  \item 单次成功提高可信度，但不能把 posterior 直接推到 1。
\end{itemize}

\teachervoice{00:53:33--00:56:04，老师把这个案例分成 multi-year human pipeline 的 recapitulation 与 B2R node 的 prospective validation。两种贡献不能混成一句“AI 发现病因”。}

\subsection{与 base LLM 的比较怎样读}

课程表格比较 AI co-scientist 与 Claude 4 / GPT-5 对若干机制要素的覆盖：初始 bradykinin accumulation、APOE4-specific B2R sensitization、neuronal energetic crisis、特异性激酶和最终因果复杂度。表格想说明多智能体系统能形成更完整的机制链，但它仍是对输出内容的分析，不是对治疗效果的随机对照实验。

\lecturefigure{slide-019.jpg}{AI co-scientist 与 base LLM 的机制要素比较}{Stanford CS25 V6 Lecture 07 官方录像 00:56:04}

\begin{knowledgebox}{读图：比较的是“机制覆盖”，不是医学能力总排名}
每一行是研究者事后关心的机制要素，列中记录系统是否命中或遗漏。AI co-scientist 的优势可能来自更长搜索、专门 Agent、工具和 tournament，而不一定来自基础模型本身更懂医学。若评价行由看到目标答案的人制定，还要警惕 rubric leakage。最稳健用途是做 ablation：逐项去掉 agent、工具或 compute，看哪些机制随之消失。
\end{knowledgebox}

\begin{warningbox}{Benchmark table 的三个限制}
样本量小；机制行可能经过人工选择；“命中一个词”与“给出可验证因果模型”不同。该表支持系统设计假设，不能独立支持疾病机制或临床决策。
\end{warningbox}

\subsection{本章小结}

Alzheimer's 案例展示了一个更完整的科研闭环：公开背景提供线索，多智能体系统重建 ACE--bradykinin--B2R 机制，合作团队再对 B2R 做前瞻干预。它比隐藏结果复现更接近发现，但课堂数据仍未发表。正确阅读方式是追踪每条因果边、冻结预测时间，并把 benchmark、实验和临床证据严格分层。

